\chapter{结果分析与讨论}

例如由于起初未能真正掌握各种控件的功能，我设想是要一个下拉菜单，但是学识肤浅的我试了很多种就是达不到我要的效果，经过查阅资料和向老师同学请教，最终解决了这个问题。

关于激光入射功率密度对导轨滚道表面硬化层深和显微硬度的影响如表 \ref{tab:4-1} 所示。下面改用可自动跨页、自动重复表头的 \texttt{longtable} 实现，后续只需要继续增减数据行即可，无需再手动拆分页和手写“续表”。

\begin{upclongtable}
  {>{\centering\arraybackslash}p{1.8cm}>{\centering\arraybackslash}p{3.8cm}>{\centering\arraybackslash}p{2.0cm}>{\centering\arraybackslash}p{3.0cm}>{\centering\arraybackslash}p{2.2cm}}
  {激光入射功率密度对导轨滚道表面硬化层深和显微硬度的影响}
  {tab:4-1}
  \toprule
  \upchdr{试验编号}{/（个）}
  & \upchdr{功率密度}{/(W$\cdot$cm$^{-2}$)}
  & \upchdr{辐照时间}{/s}
  & \upchdr{显微硬度}{/HV}
  & \upchdr{硬化层深}{/mm} \\
  \midrule
  \endfirsthead

  \upclongtablecont{5}\\
  \toprule
  \upchdr{试验编号}{/（个）}
  & \upchdr{功率密度}{/(W$\cdot$cm$^{-2}$)}
  & \upchdr{辐照时间}{/s}
  & \upchdr{显微硬度}{/HV}
  & \upchdr{硬化层深}{/mm} \\
  \midrule
  \endhead

  \bottomrule
  \endfoot

  t-1  & $6.37\times10^3$ & 0.067 & 570，456 & 0.354 \\
  t-2  & $6.37\times10^3$ & 0.067 & 570，456 & 0.354 \\
  t-3  & $6.37\times10^3$ & 0.067 & 570，456 & 0.354 \\
  t-4  & $6.37\times10^3$ & 0.067 & 570，456 & 0.354 \\
  t-5  & $6.37\times10^3$ & 0.067 & 570，456 & 0.354 \\
  t-6  & $6.37\times10^3$ & 0.067 & 570，456 & 0.354 \\
  t-7  & $6.37\times10^3$ & 0.067 & 570，456 & 0.354 \\
  t-8  & $6.37\times10^3$ & 0.067 & 570，456 & 0.354 \\
  t-9  & $6.37\times10^3$ & 0.067 & 570，456 & 0.354 \\
  t-10 & $6.37\times10^3$ & 0.067 & 570，456 & 0.354 \\
  t-11 & $6.37\times10^3$ & 0.067 & 570，456 & 0.354 \\
  t-12 & $6.37\times10^3$ & 0.067 & 570，456 & 0.354 \\
  t-13 & $6.37\times10^3$ & 0.067 & 570，456 & 0.354 \\
  t-14 & $6.37\times10^3$ & 0.067 & 570，456 & 0.354 \\
  t-15 & $6.37\times10^3$ & 0.067 & 570，456 & 0.354 \\
  t-16 & $6.37\times10^3$ & 0.067 & 570，456 & 0.354 \\
\end{upclongtable}
